\chapter{Cross-chapter statistical conventions}
\label{app:shared-conventions}

This appendix is a concordance, not a replacement for the model definitions in
Chapters~\ref{ch:research-a}--\ref{ch:research-d}.  The four projects share
ideas and computational devices, but they do not always use the same latent
variable scaling, inferential target, scoring rule, or meaning of an interval.
Keeping those distinctions explicit prevents an algebraic resemblance from
being read as scientific equivalence.

\section{Targets and update mechanisms}

Let \(p\in(0,1)\) denote a response-quantile level.  In
Chapters~\ref{ch:research-a}--\ref{ch:research-c}, a fitted location is intended
to represent a conditional \(p\)-quantile under an AL or exAL working
likelihood.  A posterior or variational interval for that fitted location is a
parameter- or state-uncertainty summary.  A response-level predictive interval
also includes the observation component of the working model.  Neither object
is, without a separate construction, a frequentist tolerance interval.

Chapter~\ref{ch:research-d} instead begins with a fixed-content interval as the
estimand.  Its content \(c\) determines the population window, its tilt
\(\delta\) determines placement through a retained-mean condition, and its
generalized-Bayes learning rate determines the concentration of a loss-based
update.  These three quantities are not interchangeable.  In particular, the
pseudo-AL identity used for computation in that chapter is an augmentation of
the scalar product loss; it is not an AL response likelihood for either
endpoint.

\section{AL and exAL latent-variable scales}

Chapters~\ref{ch:research-a} and~\ref{ch:research-c} use the convention

\begin{align*}
v_t\mid\sigma &\sim \operatorname{Exp}(\text{rate}=1/\sigma),\\
Y_t\mid v_t,\ldots &\sim
\mathcal N\{\mu_t+A_pv_t,\;\sigma B_pv_t\},
\end{align*}

with the corresponding exAL shift and truncated-normal latent variable added
when required.  Thus \(v_t\) has mean \(\sigma\).  Chapter~\ref{ch:research-b}
uses a standardized exponential latent variable in parts of its source-aware
formulation.  Writing \(v_t=\sigma z_t\), with
\(z_t\sim\operatorname{Exp}(1)\), converts between the two representations and
changes the displayed Gaussian variance from \(\sigma B_pv_t\) to
\(B_p\sigma^2z_t\).  Posterior formulas must therefore be read under the local
convention; factors of \(\sigma\) cannot be transferred between chapters by
notation alone.

The normal distribution is parameterized by mean and variance throughout the
dissertation.  Inverse-gamma and generalized inverse Gaussian conventions are
stated in the technical appendix in which they are used.  The symbol
\(\sigma\) is consequently a model-specific scale, not a promise that all four
chapters assign it the same likelihood interpretation or prior.

\section{Multiple quantiles and predictive distributions}

Separately fitted conditional quantiles, a stacked composite working
likelihood, and a predictive distribution are distinct objects.  The
\pkg{exdqlm} synthesis in Chapter~\ref{ch:research-a} interpolates and, when
requested, monotonizes posterior-predictive information from separate fits.
The hydrologic procedure in Chapter~\ref{ch:research-b} synthesizes
source-aware predictive information across levels.  Chapter~\ref{ch:research-c}
also studies a stacked multi-quantile readout and reports crossing before any
monotone rearrangement.  None of these constructions is described as a unique
joint distribution unless the chapter supplies the additional distributional
assumptions needed for that statement.

\section{Scores and empirical comparisons}

The continuous ranked probability score (CRPS) is used only when a predictive
distribution or response-level sample is available.  The hydrologic chapter
reports its empirical predictive CRPS under the source-aware forecast design.
The Q--DESN chapter also uses a finite-grid integrated quantile score, denoted
there by \(\operatorname{aCRPS}\), over the stated fitted probability range.
That finite-grid quantity depends on the grid and tail convention and must not
be compared numerically with a full-distribution CRPS from another study.
Check loss evaluates a specified conditional quantile.  Interval scores and
tolerance calibration evaluate still different actions.  Every empirical
comparison is therefore interpreted within its recorded target, horizon,
data split, and score definition.

\section{Uncertainty language}

The dissertation uses \emph{credible interval} for posterior uncertainty in a
parameter, latent state, or fitted functional; \emph{predictive interval} or
\emph{predictive band} for a future response under the stated model; and
\emph{tolerance interval} for a procedure calibrated to cover a specified
population proportion with a specified repeated-sampling confidence.  An
empirical coverage frequency is evidence about a particular validation design,
not a proof of a general coverage guarantee.  Likewise, a variational
approximation is assessed as an approximation to its declared posterior
target; computational speed does not strengthen the associated scientific
claim.

\section{Where the complete details reside}

Appendix~\ref{app:exdqlm-technical} gives the exDQLM posterior and software
calculations; Appendix~\ref{app:hydrology-technical} gives the hydrologic MCMC,
VB, and secondary diagnostics; Appendix~\ref{app:qdesn-technical} gives the
Q--DESN prior, posterior, forecasting, and sensitivity calculations; and
Appendix~\ref{app:mti-technical} gives the interval proofs, calibration
protocol, and proposed extension derivations.  Those local appendices are
authoritative whenever this concordance deliberately suppresses notation.
